\documentclass[11pt]{article}

\usepackage[utf8]{inputenc}
\usepackage[T1]{fontenc}
\usepackage{geometry}
\usepackage{amsmath}
\usepackage{amssymb}
\usepackage{booktabs}
\usepackage{hyperref}
\usepackage{xurl}
\usepackage{xcolor}
\usepackage{listings}
\usepackage{enumitem}
\usepackage{parskip}
\usepackage{graphicx}
\graphicspath{{../images/}}

\geometry{margin=1in}

\lstset{
  language=Java,
  basicstyle=\ttfamily\small,
  keywordstyle=\color{blue},
  commentstyle=\color{gray},
  stringstyle=\color{teal},
  showstringspaces=false,
  columns=fullflexible,
  frame=single,
  framerule=0.2pt,
  breaklines=true
}

\setlist[itemize]{topsep=0.3em,itemsep=0.2em}
\setlist[enumerate]{topsep=0.3em,itemsep=0.2em}

% Simple per-section source line.
\newcommand{\notesources}[1]{\par\noindent{\small\color{gray}\textit{Sources:} \sloppy #1\par}}

% Unit-wise source strings for this subject (used by slides/notes).
% Path is relative to the lecture `latex/` folder (the compilation CWD).
% OOPS with Java (CSEG1044) — unit-wise sources
%
% These are short, student-facing reference strings intended to be used with:
%   - \slidesources{...}  (in beamer slides)
%   - \notesources{...}   (in lecture notes)
%
% Style inspiration: other subjects in My-Subjects/ use semicolon-separated sources
% and include an "accessed" date for web documentation.

\newcommand{\OOPSUnitISources}{Schildt, \textit{Java: A Beginner's Guide} (10e), McGraw-Hill (Java fundamentals + OOP basics).; Oracle Java Tutorials: Getting Started + Language Basics + Classes and Objects (accessed \today).; Java SE API docs (e.g., \texttt{String}, \texttt{Scanner}) (accessed \today).}

\newcommand{\OOPSUnitIISources}{Schildt, \textit{Java: The Complete Reference} (12e), McGraw-Hill (classes, inheritance, packages, interfaces).; Bloch, \textit{Effective Java} (3e), Addison-Wesley, 2018 (design choices; interfaces vs abstract classes).; Oracle Java Tutorials: Classes and Objects + Interfaces + Packages (accessed \today).; Java SE API docs (e.g., \texttt{Comparable}, \texttt{Runnable}) (accessed \today).}

\newcommand{\OOPSUnitIIISources}{Schildt, \textit{Java: The Complete Reference} (12e) (nested/inner classes, exceptions, threads, I/O).; Oracle Java Tutorials: Nested Classes + Exceptions + Concurrency + Basic I/O (accessed \today).; Java SE API docs (e.g., \texttt{Thread}, \texttt{AutoCloseable}) (accessed \today).}

\newcommand{\OOPSUnitIVSources}{Schildt, \textit{Java: The Complete Reference} (12e) (generics, lambdas, Swing, JDBC).; Bloch, \textit{Effective Java} (3e) (generics + lambdas).; Oracle Java Tutorials: Generics + Lambda Expressions + Swing + JDBC Basics (accessed \today).; Java SE API docs (e.g., \texttt{java.util.function}, \texttt{javax.swing}, \texttt{java.sql}) (accessed \today).}

\newcommand{\OOPSUnitVSources}{Schildt, \textit{Java: The Complete Reference} (12e) (Collections Framework + wrapper classes).; Oracle Java docs: Collections Framework overview (accessed \today).; Java SE API docs (\texttt{java.util} collections; wrapper classes in \texttt{java.lang}) (accessed \today).}

\newcommand{\OOPSUnitVISources}{Git SCM: \textit{Pro Git} (2e) (version control workflow), accessed \today.; GitHub Docs (repositories, issues, pull requests), accessed \today.; Oracle Java Tutorials: Swing + JDBC (accessed \today).; JUnit 5 User Guide (accessed \today).; Apache Maven Project: Getting Started (accessed \today).; Gradle User Manual: Getting Started (accessed \today).; UML class diagrams: UML 2.x overview/spec (accessed \today).}

\newcommand{\OOPSMiscWhyInterfacesSources}{Schildt, \textit{Java: The Complete Reference} (12e) (interfaces).; Bloch, \textit{Effective Java} (3e) (prefer interfaces where appropriate).; Oracle Java Tutorials: Interfaces (accessed \today).; Java SE API docs (e.g., \texttt{Comparable}, \texttt{Runnable}) (accessed \today).}



\title{Object Oriented Programming Lab (CSEG1044)\\Experiment-03\&04: Operators + Control Flow (Student Handout --- Before Submission)}
\author{Tofik Ali}
\date{\today}

\begin{document}
\maketitle
\tableofcontents

\section{About this handout}
This is a \textbf{pre-submission} handout.
Do the implementation yourself; this document provides the required tasks and some hints.

\section{Experiment 03 (Syllabus)}
\begin{itemize}[leftmargin=*]
  \item \textbf{Objective:} To apply data types and operators in Java programs.
  \item \textbf{Problem Statement:} Write a Java program demonstrating the use of primitive data types and operators.
  \item \textbf{Expected Outcome:} Students will understand data manipulation using operators.
\end{itemize}

\section{Experiment 04 (Syllabus)}
\begin{itemize}[leftmargin=*]
  \item \textbf{Objective:} To implement control flow using conditional and looping statements.
  \item \textbf{Problem Statement:} Design a Java program using decision-making and looping constructs.
  \item \textbf{Expected Outcome:} Students will be able to control program flow effectively.
\end{itemize}

\section{What you must do (minimum requirements)}

\subsection{Task A --- Types + operators mini-lab (Exp-03)}
Create \texttt{TypesOperatorsDemo.java} that:
\begin{enumerate}[leftmargin=*]
  \item reads two integers \texttt{a} and \texttt{b},
  \item prints \texttt{a+b}, \texttt{a-b}, \texttt{a*b},
  \item if \texttt{b != 0}, prints:
  \begin{itemize}[leftmargin=*]
    \item integer division: \texttt{a / b}
    \item floating division: \texttt{a / (double) b}
    \item remainder: \texttt{a \% b}
  \end{itemize}
  \item demonstrates precedence:
  \begin{itemize}[leftmargin=*]
    \item \texttt{2 + 3 * 4}
    \item \texttt{(2 + 3) * 4}
  \end{itemize}
  \item demonstrates casting:
  \begin{itemize}[leftmargin=*]
    \item widening: \texttt{int -> double}
    \item narrowing: \texttt{double -> int} (show truncation)
  \end{itemize}
  \item prints one boolean expression using \texttt{\&\&} and \texttt{||}.
\end{enumerate}

\subsection{Task B --- Menu-driven utility program (Exp-04)}
Create \texttt{MenuUtility.java} that repeats until Exit:
\begin{enumerate}[leftmargin=*]
  \item Even/Odd check
  \item Factorial (loop)
  \item Prime check
  \item Exit
\end{enumerate}

\noindent Requirements:
\begin{itemize}[leftmargin=*]
  \item Use \texttt{switch} for menu choice handling.
  \item Use loops for factorial and prime check.
  \item Handle invalid menu choice.
  \item Show at least one \texttt{continue} usage (example: skip processing if input is negative).
\end{itemize}

\section{Hints (read only if stuck)}
\begin{itemize}[leftmargin=*]
  \item \texttt{5/2} is integer division (result \texttt{2}); use \texttt{5/2.0} for \texttt{2.5}.
  \item Always check \texttt{b==0} before division/modulo.
  \item For prime check, loop only up to \texttt{i*i <= n}.
  \item In \texttt{switch}, forgetting \texttt{break} causes fall-through.
\end{itemize}

\section{What to submit (MANDATORY)}
Submit a single zipped folder containing:
\begin{itemize}[leftmargin=*]
  \item \texttt{TypesOperatorsDemo.java}
  \item \texttt{MenuUtility.java}
  \item sample output for both programs
  \item observations (5--10 lines): integer vs floating division; where loops are used; common mistakes you faced
\end{itemize}

\notesources{\OOPSUnitISources}
\end{document}
